\documentclass[10pt,a4paper]{amsart}
\usepackage[margin=19mm]{geometry}
\usepackage{amsmath,amssymb,mathtools}
\usepackage[scr=rsfs,cal=euler]{mathalfa}
\usepackage{xfrac}
\usepackage[unicode,hidelinks,bookmarks=false]{hyperref}
\newcommand{\ie}{i.e.\ }
\newcommand{\cf}{cf.\ }
\newcommand{\resp}{respectively\ }
\newcommand{\Subsection}{\subsection}
\providecommand{\nobreakdash}{\nobreak\hbox{-}}
\setlength{\emergencystretch}{2em}
\multlinegap0pt
\usepackage{mathtools}
\usepackage{amssymb}
\usepackage[shortlabels]{enumitem}
\usepackage{float}
\usepackage{cancel}
\newtheorem{lemma}{Lemma}
\newcommand{\R}{\mathbb{R}}
\newcommand{\eps}{\varepsilon}
\title{A Phase-Field Method for Curvature Flow of Networks\\ with Triple Junction Drag}
\author{Yuchuan Yang \and Selim Esedo\={g}lu}
\date{Source: arXiv:2608.00293v1 (CC BY 4.0)}
\begin{document}
\maketitle

\noindent\emph{Source and licence.} A Phase-Field Method for Curvature Flow of Networks\\ with Triple Junction Drag. Version arXiv:2608.00293v1 (CC BY 4.0), distributed by arXiv under the Creative Commons Attribution 4.0 International licence, http://creativecommons.org/licenses/by/4.0/, as recorded in the arXiv licence metadata for this exact version and shown on the abstract page https://arxiv.org/abs/2608.00293v1. Full licence notice: \texttt{licenses/arxiv-2608.00293-CC-BY-4.0.txt}.\par\smallskip

\noindent\textit{Editorial scope.} Counted unit: the article's lemma lem:asymptotic (source0.tex lines 670--675). The article's complete original proof of it is delivered with it (source0.tex lines 677--729, one \texttt{proof} environment of 53 lines). No mathematics is written, repaired or re-derived: the statement and the proof are the article's own lines, in the article's own notation, and no retained line is substituted or re-flowed.  No further source-local context is needed: every cross-reference of every retained line resolves inside the two delivered spans.  Nothing was omitted from any retained span, so the package carries no omission record.  No retained line contains a citation, so the wrapper carries no bibliography and no bibliography entry.  No retained line contains a figure environment, an image-inclusion command or drawing code, so no image is delivered and the package declares no asset dependency.  Editorial formatting only, in addition: an AMS \texttt{amsart} preamble that restates the article's own declarations and macros used by the retained lines, namely the environments \texttt{lemma} on separate counters, together with the standard packages those lines need.  The retained environments print in this document as Lemma 1 in order of appearance, whereas the article prints the same items with the numbering of its own full document.  The complete native source of the article as distributed in the arXiv e-print for this exact version is bundled unchanged as \texttt{sources/206532/source0.tex}.
\begin{lemma}\label{lem:asymptotic}
    \begin{equation}
        J_{u^0} B (B + \eps^2 I_{2\times 2})^{-2}B = J_{u^0}I_{2\times 2} + \mathcal{O}(\eps)
    \end{equation}
    where the constant appearing in the error term depends only on the largest singular value of $\nabla u^0$.
\end{lemma}

\begin{proof}
    We first orthogonally diagonalize 
    \begin{equation}
        B = Q \Lambda Q^\top, \qquad \Lambda = \begin{pmatrix}
            \lambda_1 & 0 \\ 0 & \lambda_2
        \end{pmatrix},
        \qquad \lambda_1,\lambda_2 \geq 0
    \end{equation}
    and note that $J_{u^0} = \sqrt{\lambda_1 \lambda_2}$. Then,
    \begin{equation}
        J_{u^0} B (B + \eps^2 I_{2\times 2})^{-2}B = Q\begin{pmatrix}
            \mu_1(\eps) & 0 \\ 0 & \mu_2(\eps)
        \end{pmatrix}
        Q^\top
    \end{equation}
    where 
    \begin{equation}
        \mu_i(\eps) \coloneqq \frac{\lambda_i^2 \sqrt{\lambda_1 \lambda_2}}{(\lambda_i + \eps^2)^2}, \qquad i = 1,2.
    \end{equation}
    Our goal is to show that 
    \begin{equation}
        |\mu_i(\eps) - \sqrt{\lambda_1 \lambda_2}| \leq C\eps
    \end{equation}
    for some constant $C$. WLOG, we can just consider the case when $i=1$.
    \begin{equation}
        \begin{split}
            |\mu_i(\eps) - \sqrt{\lambda_1 \lambda_2}| &= \sqrt{\lambda_1 \lambda_2}\bigg|1-\frac{\lambda_1^2}{(\lambda_1 + \eps^2)^2}\bigg|\\
            %%
            &= \eps\sqrt{\lambda_2}f(\eps),
        \end{split}
    \end{equation}
    where 
    \begin{equation}
        f(\eps) \coloneqq \sqrt{\lambda_1}\bigg(\frac{2\lambda_1 \eps + \eps^3}{(\lambda_1 + \eps^2)^2}\bigg)\geq 0.
    \end{equation}
    We claim that 
    \begin{equation}
        \sup_{\eps>0} f(\eps) \leq 1.
    \end{equation}
    Indeed, setting $z = \eps/\sqrt{\lambda_1}$, we get
    \begin{equation}
        f(\eps) = \frac{z(2+z^2)}{(1+z^2)^2} \leq \frac{2z}{1+z^2} \leq 1.
    \end{equation}
    Thus,
    \begin{equation}
        |\mu_i(\eps) - \sqrt{\lambda_1 \lambda_2}| \leq \sqrt{\lambda_2}\eps \leq C\eps
    \end{equation}
    where 
    \begin{equation}
        C \coloneqq \sup_{y\in\R^2} \sqrt{\lambda_{\max}((\nabla u^0)^\top \nabla u^0)}.
    \end{equation}
    
\end{proof}

\end{document}
